\section{Discussion}
\label{sec:discussion}

The results provide nuanced support for our hypothesis that 'mode' acts as a structured latent variable in text interpretation.

\subsection{Interpretation of Results}

{\bf Detectability vs. Artifact Granularity.}
The contrast between the high recall for \aigen text and the total failure for \dictated text suggests a hierarchy of artifact saliency.
LLMs appear highly sensitive to the "cleanliness" and stylistic consistency of AI-generated text.
However, they struggle to zero-shot distinguish between different types of human errors.
Phonetic homophones (e.g., "there/their"), while systematic, are likely 'corrected' by the model's internal language prior during the reasoning process, leading it to misidentify \dictated text as either clean (\aigen) or orthographically noisy (\typed).

{\bf Latent Separation.}
The Silhouette Score of \silhouette and the high variance explained by PC1 (\pcaone) confirm that origin mode is not just a surface artifact but a dimension of variance captured by modern embeddings.
The primary axis of separation appears to be a "Clean vs. Noisy" dimension.
The fact that \dictated and \typed modes overlap suggests that separating orthographic from phonetic noise requires more specialized representations or few-shot demonstration to calibrate the model's sensitivity.

\subsection{Limitations}

{\bf Synthetic Error Injection.}
Our reliance on synthetic homophone and typo injection, while necessary for content control, may introduce artifacts that differ from organic human data.
Real-world dictation errors often involve more complex ASR artifacts beyond simple homophone substitution.

{\bf Model Scale.}
The use of \gptmini may have limited the zero-shot reasoning capabilities.
Larger models like GPT-4o or Claude 3.5 Sonnet might possess the granularity required to detect subtle phonetic artifacts that a smaller model glosses over.

{\bf Sample Size.}
With 150 total samples, this study acts as a proof-of-concept.
A larger-scale evaluation across diverse topics and languages is needed to confirm the generalizability of these latent modes.

\subsection{Broader Implications}

The ability to isolate 'mode' as a latent variable has significant implications for forensic linguistics and data curation.
By identifying the origin mode, systems can adapt their error correction strategies—for instance, using phonetic models for \dictated text and orthographic models for \typed text.
Furthermore, the isolation of \aigen signatures is vital for detecting model-in-the-loop content in sensitive domains like legal or academic writing.
